\documentclass[12pt]{article}

\usepackage[a4paper,left=20mm,right=25mm,top=25mm,bottom=25mm]{geometry}
\usepackage[utf8]{inputenc}
\usepackage[ngerman]{babel}
\usepackage[T1]{fontenc}

% Tabellen
\usepackage{array}
\usepackage{tabularx}
\usepackage{booktabs}
\newcolumntype{L}[1]{>{\raggedright\arraybackslash}p{#1}}

% Matheformeln
\usepackage{amsmath,amsfonts,amssymb,amsthm}
\newtheorem{thm}{Satz}
\newtheorem{defn}[thm]{Definition}
\newtheorem{lem}[thm]{Lemma}
\newtheorem{prop}[thm]{Proposition}
\newtheorem{cor}[thm]{Korollar}

\usepackage{hyperref}
\usepackage{mathptmx}
\usepackage[scaled=.90]{helvet}

% **********************************************************************
\begin{document}

\section{Das Problem}

Es sei $n\in\mathbb{N}$ mit $n\geq 2$. Wir definieren die Indexmenge
$N=\{1,2,\dots,n\}$ und zwei Vektoren
\[
p=(p_1,p_2,\dots,p_n),
\qquad
q=(q_1,q_2,\dots,q_n)
\]
mit $p_i,q_i\in\{0,1\}$ für $i\in N$. Für $i,j\in N$ mit $i<j$ sei
\begin{align*}
\mathrm{one}_i^j(p)
&=\#\{k\mid i<k<j,\ p_k=1\},\\
\overline{\mathrm{one}}_i^j(p)
&=\#\{k\mid 1\leq k<i,\ p_k=1\}
 +\#\{k\mid j<k\leq n,\ p_k=1\},\\
\mathrm{zero}_i^j(p)
&=\#\{k\mid i<k<j,\ p_k=0\},\\
\overline{\mathrm{zero}}_i^j(p)
&=\#\{k\mid 1\leq k<i,\ p_k=0\}
 +\#\{k\mid j<k\leq n,\ p_k=0\}.
\end{align*}

\begin{defn}
Es sei das Tupel $(n,p,q)$ wie oben beschrieben. Gibt es $i,j\in N$ mit
$i<j$, $p_i=p_j$ und
\begin{align}
\mathrm{one}_i^j(p)
&>\overline{\mathrm{one}}_i^j(p),
\label{eq:ones-greater-original-notation}\\
\mathrm{zero}_i^j(p)
&>\overline{\mathrm{zero}}_i^j(p),
\label{eq:zeros-greater-original-notation}
\end{align}
dann kann für $(i,j)$ folgende Transformation durchgeführt werden, die wir
Tausch oder Wechsel nennen:
\[
\tau_{ij}(p)
=(p_1,\dots,p_{i-1},1-p_i,p_{i+1},\dots,
  p_{j-1},1-p_j,p_{j+1},\dots,p_n).
\]
Kann $q$ durch eine endliche Folge von $m$ Wechseln aus $p$ erreicht werden,
also
\[
(\tau_{i_mj_m}\circ\cdots\circ\tau_{i_2j_2}
 \circ\tau_{i_1j_1})(p)=q,
\]
so schreiben wir $p\sim q$.
\end{defn}

\begin{defn}
Es sei $(n,p,q)$ gegeben mit
\[
\#\{i\in N\mid p_i=1\}
\quad\text{und}\quad
\#\{i\in N\mid p_i=0\}
\]
ungerade. Gilt $p\sim q$, so nennen wir $(n,p,q)$ lösbar.
\label{defn:problembeschreibung-median}
\end{defn}

Die zu beantwortende Frage lautet: Ist das gegebene Tupel $(n,p,q)$ lösbar?

\section{Grundlegende Eigenschaften}

\begin{prop}
Es sei $(n,p,q)$ wie in Definition
\ref{defn:problembeschreibung-median} gegeben. Dann ist $n$ gerade.
\end{prop}

\begin{proof}
Die Anzahl der Einsen und die Anzahl der Nullen in $p$ sind ungerade. Ihre
Summe ist $n$ und damit gerade.
\end{proof}

\begin{prop}[Paritätsinvarianz]
Bei jedem Wechsel ändern sich die Anzahl der Einsen und die Anzahl der Nullen
jeweils um zwei. Insbesondere bleiben beide Anzahlen ungerade.
\label{prop:paritaet-median}
\end{prop}

\begin{proof}
Da $p_i=p_j$ gilt, werden entweder zwei Nullen zu Einsen oder zwei Einsen zu
Nullen.
\end{proof}

\begin{prop}[Umkehrbarkeit]
Kann $(i,j)$ in $p$ wechseln und ist $p'=\tau_{ij}(p)$, so kann $(i,j)$ auch
in $p'$ wechseln. Insbesondere ist die Relation $\sim$ symmetrisch.
\label{prop:umkehrbar-median}
\end{prop}

\begin{proof}
Die Spieler $i$ und $j$ stimmen weder vor noch nach dem Wechsel ab. Alle
anderen Einträge und damit alle Ja- und Nein-Stimmen bleiben unverändert.
Die Ungleichungen
\eqref{eq:ones-greater-original-notation} und
\eqref{eq:zeros-greater-original-notation} gelten deshalb auch für den
Rückwechsel.
\end{proof}

\section{Mediane und Medianintervall}

\begin{defn}
Es sei $s\in\{0,1\}^n$, und die Anzahl der Nullen sowie die Anzahl der Einsen
in $s$ seien ungerade. Für $c\in\{0,1\}$ definieren wir
\[
P_c(s)=\{i\in N\mid s_i=c\}.
\]
Ist $\#P_c(s)=2h_c+1$, so bezeichnen wir die Position des $(h_c+1)$-ten
Vorkommens von $c$ mit $m_c(s)$ und nennen sie den Median der $c$-Spieler.
Weiter definieren wir
\begin{align*}
\ell(s)&=\min\{m_0(s),m_1(s)\},\\
r(s)&=\max\{m_0(s),m_1(s)\}
\end{align*}
und die Mediansignatur
\[
M(s)=\bigl(\ell(s),r(s),s_{\ell(s)}s_{\ell(s)+1}\cdots s_{r(s)}\bigr).
\]
Das Teilwort $s_{\ell(s)}\cdots s_{r(s)}$ nennen wir Medianintervall.
\end{defn}

\begin{lem}
Es sei $s\in\{0,1\}^n$ mit einer ungeraden Anzahl Nullen und Einsen. Können
$i<j$ wechseln, so gilt
\[
i<m_0(s)<j
\qquad\text{und}\qquad
i<m_1(s)<j.
\]
\label{lem:beide-mediane-dazwischen}
\end{lem}

\begin{proof}
Wir betrachten zunächst den Wert $c=s_i=s_j$. Es gebe $2h+1$ Einträge mit
Wert $c$, und $i$ beziehungsweise $j$ sei das $a$-te beziehungsweise $b$-te
Vorkommen von $c$. Die Anzahl der Ja-Stimmen der $c$-Spieler ist dann
$b-a-1$, die Anzahl ihrer Nein-Stimmen ist
\[
(a-1)+(2h+1-b).
\]
Die strikte Mehrheitsbedingung liefert
\[
b-a-1>(a-1)+(2h+1-b)
\]
und damit $b-a\geq h+1$. Wäre $a\geq h+1$ oder $b\leq h+1$, so wäre wegen
$1\leq a<b\leq 2h+1$ dagegen $b-a\leq h$. Es folgt
\[
a<h+1<b,
\]
also liegt $m_c(s)$ strikt zwischen $i$ und $j$.

Für den anderen Wert $1-c$ gebe es $2k+1$ Einträge. Keiner von ihnen ist ein
Endpunkt. Die zugehörige strikte Ja-Mehrheit bedeutet, dass mindestens $k+1$
dieser Einträge zwischen $i$ und $j$ liegen. Diese Einträge bilden einen
zusammenhängenden Block in der geordneten Liste aller $(1-c)$-Vorkommen. Jeder
solche Block enthält das $(k+1)$-te Vorkommen. Somit liegt auch $m_{1-c}(s)$
strikt zwischen $i$ und $j$.
\end{proof}

\begin{lem}[Invarianz der Mediansignatur]
Können $(i,j)$ in $s$ wechseln und ist $s'=\tau_{ij}(s)$, so gilt
\[
M(s')=M(s).
\]
\label{lem:mediansignatur-invariant}
\end{lem}

\begin{proof}
Nach Lemma~\ref{lem:beide-mediane-dazwischen} liegen beide Mediane strikt
zwischen $i$ und $j$. Für den Wert $s_i=s_j$ verschwindet somit je ein
Vorkommen links und rechts von seinem Median. Für den Wert $1-s_i$ kommt je
ein Vorkommen links und rechts von seinem Median hinzu. Beide bisherigen
Mediane besitzen danach weiterhin gleich viele gleichartige Einträge auf
beiden Seiten und bleiben daher unverändert.

Außerdem gilt
\[
i<\ell(s)\leq r(s)<j.
\]
Kein Eintrag des Medianintervalls wird umgeschaltet. Seine Grenzen und sein
gesamtes Teilwort bleiben gleich.
\end{proof}

\begin{cor}
Ist $(n,p,q)$ lösbar, so gilt $M(p)=M(q)$.
\label{cor:mediansignatur-notwendig}
\end{cor}

\section{Wechsel außerhalb des Medianintervalls}

\begin{lem}[Wechsel mit den Randpositionen]
Es sei $s\in\{0,1\}^n$ wie oben und $[\ell(s),r(s)]$ sein Medianintervall.
Dann gilt:
\begin{enumerate}
\item Ist $i<\ell(s)$ und $s_i=s_n$, so können $(i,n)$ wechseln.
\item Ist $j>r(s)$ und $s_1=s_j$, so können $(1,j)$ wechseln.
\end{enumerate}
\label{lem:wechsel-mit-rand}
\end{lem}

\begin{proof}
Wir zeigen die erste Aussage; die zweite folgt symmetrisch. Es sei
$c=s_i=s_n$, und es gebe $2h+1$ Einträge mit Wert $c$. Weil $i$ vor beiden
Medianen liegt, gibt es vor $i$ höchstens $h-1$ weitere $c$-Einträge. Daher
stimmen höchstens $h-1$ der verbleibenden $2h-1$ Spieler mit Wert $c$ mit
Nein und mindestens $h$ mit Ja.

Es gebe $2k+1$ Einträge mit Wert $1-c$. Da auch deren Median rechts von $i$
liegt, befinden sich höchstens $k$ von ihnen vor $i$. Folglich stimmen
höchstens $k$ mit Nein und mindestens $k+1$ mit Ja. Damit gelten je nach Wert
von $c$ sowohl Ungleichung
\eqref{eq:ones-greater-original-notation} als auch Ungleichung
\eqref{eq:zeros-greater-original-notation}.
\end{proof}

\begin{lem}[Randkonstruktion]
Es seien $i<\ell(s)$ und $j>r(s)$ mit $s_i=s_j$. Dann gibt es eine Folge von
höchstens drei Wechseln, deren Gesamteffekt genau dem Umschalten von $s_i$ und
$s_j$ entspricht. Alle anderen Einträge besitzen am Ende wieder ihren
ursprünglichen Wert.
\label{lem:randkonstruktion-original}
\end{lem}

\begin{proof}
Für $i=1$ oder $j=n$ ist der direkte Wechsel nach
Lemma~\ref{lem:wechsel-mit-rand} möglich. Es sei daher
\[
1<i<\ell(s)\leq r(s)<j<n
\]
und $x=s_i=s_j$. Abhängig von den beiden Randwerten führen wir die folgende
Folge von links nach rechts aus:
\begin{center}
\begin{tabularx}{\textwidth}{@{}>{\raggedright\arraybackslash}X
                                  >{\raggedright\arraybackslash}X@{}}
\toprule
Voraussetzung & Folge der Wechsel \\
\midrule
$s_1=s_n=x$
  & $\tau_{in},\ \tau_{1j},\ \tau_{1n}$ \\
$s_1=s_n\neq x$
  & $\tau_{1n},\ \tau_{in},\ \tau_{1j}$ \\
$s_1\neq s_n$ und $s_n=x$
  & $\tau_{in},\ \tau_{1n},\ \tau_{1j}$ \\
$s_1\neq s_n$ und $s_1=x$
  & $\tau_{1j},\ \tau_{1n},\ \tau_{in}$ \\
\bottomrule
\end{tabularx}
\end{center}
Bei jedem Schritt sind die beiden betroffenen Einträge gleich. Daher ist
jeder Schritt nach Lemma~\ref{lem:wechsel-mit-rand} erlaubt. Mit
Lemma~\ref{lem:mediansignatur-invariant} bleibt das Medianintervall in allen
Zwischenzuständen unverändert. In jedem der vier Fälle werden $s_1$ und $s_n$
am Ende wiederhergestellt und genau $s_i,s_j$ umgeschaltet.
\end{proof}

\begin{lem}[Bilanzlemma]
Es seien $I$ und $J$ zwei nichtleere, disjunkte Indexmengen. Eine abstrakte
Operation schalte je einen Eintrag aus $I$ und $J$ um, sofern beide aktuell
gleich sind. Zwei binäre Belegungen $u$ und $v$ lassen sich genau dann durch
solche Operationen ineinander überführen, wenn
\begin{align}
\#\{i\in I\mid u_i=1\}-\#\{j\in J\mid u_j=1\}
&=
\#\{i\in I\mid v_i=1\}-\#\{j\in J\mid v_j=1\}.
\label{eq:bilanz-original-notation}
\end{align}
\label{lem:binaere-bilanz-original}
\end{lem}

\begin{proof}
Jede Operation erhöht beide Einsenanzahlen um eins oder verringert beide um
eins. Die Differenz in Gleichung
\eqref{eq:bilanz-original-notation} ist daher invariant.

Für die umgekehrte Richtung vergleichen wir die aktuelle Belegung mit $v$.
Gibt es auf verschiedenen Seiten zwei falsche Einträge mit gleichem aktuellen
Wert, so können wir beide unmittelbar umschalten und damit korrigieren.

Gibt es auf derselben Seite zwei falsche Einträge mit verschiedenen aktuellen
Werten, wählen wir einen beliebigen Eintrag der anderen Seite als
Zwischenspeicher. Zunächst schalten wir den falschen Eintrag um, dessen Wert
dem Speicherwert entspricht. Anschließend stimmt der neue Speicherwert mit
dem zweiten falschen Eintrag überein. Ein zweiter Schritt korrigiert diesen
und stellt gleichzeitig den Speicher wieder her.

Wäre keine der beiden Reduktionen möglich, obwohl noch falsche Einträge
existieren, dann hätten alle falschen Einträge auf jeder Seite jeweils
denselben aktuellen Wert und die Werte der beiden Seiten wären verschieden.
Die Abweichungen der beiden Seiten in
Gleichung~\eqref{eq:bilanz-original-notation} hätten dann entgegengesetzte,
von null verschiedene Vorzeichen. Dies ist unmöglich. Somit kann die Anzahl
der falschen Einträge schrittweise bis auf null reduziert werden.
\end{proof}

\section{Lösung}

\begin{thm}
Es seien $p,q\in\{0,1\}^n$, wobei sowohl in $p$ als auch in $q$ die Anzahl
der Nullen und die Anzahl der Einsen ungerade sei. Dann gilt
\[
p\sim q
\quad\Longleftrightarrow\quad
M(p)=M(q).
\]
\label{thm:mediansignatur-charakterisierung}
\end{thm}

\begin{proof}
Die Richtung ``$\Rightarrow$'' folgt aus
Korollar~\ref{cor:mediansignatur-notwendig}.

Für die Richtung ``$\Leftarrow$'' gelte $M(p)=M(q)$. Wir schreiben
\[
\ell=\ell(p)=\ell(q),
\qquad
r=r(p)=r(q)
\]
und bezeichnen das gemeinsame Medianintervall mit
\[
C=p_\ell p_{\ell+1}\cdots p_r
 =q_\ell q_{\ell+1}\cdots q_r.
\]
Außerdem setzen wir
\[
I=\{1,\dots,\ell-1\},
\qquad
J=\{r+1,\dots,n\}.
\]

Zunächst zeigen wir Gleichung
\eqref{eq:bilanz-original-notation} für die Einschränkungen von $p$ und $q$
auf $I\cup J$. Liegt der Median der Einsen bei $\ell$, so besitzt jeder
Vektor $s$ mit Signatur $M(p)$ gleich viele Einsen links und rechts von
$\ell$. Daher gilt
\begin{align*}
\#\{i\in I\mid s_i=1\}
&=
\#\{k\in\{\ell,\dots,r\}\mid s_k=1\}-1
 +\#\{j\in J\mid s_j=1\}.
\end{align*}
Somit ist
\begin{align*}
\#\{i\in I\mid s_i=1\}
-\#\{j\in J\mid s_j=1\}
&=
\#\{k\in\{\ell,\dots,r\}\mid s_k=1\}-1,
\end{align*}
also ausschließlich durch $C$ bestimmt. Liegt der Median der Einsen bei $r$,
so folgt analog
\begin{align*}
\#\{i\in I\mid s_i=1\}
-\#\{j\in J\mid s_j=1\}
&=
1-\#\{k\in\{\ell,\dots,r\}\mid s_k=1\}.
\end{align*}
Die Bilanz ist daher in $p$ und $q$ gleich.

Sind $I$ und $J$ nichtleer, so liefert
Lemma~\ref{lem:binaere-bilanz-original} eine Folge abstrakter Umschaltungen,
welche $p$ außerhalb von $C$ in $q$ überführt. Jede abstrakte Umschaltung
zweier gleicher Einträge $p_i=p_j$ mit $i\in I$ und $j\in J$ kann nach
Lemma~\ref{lem:randkonstruktion-original} durch höchstens drei erlaubte
Wechsel realisiert werden. Das Medianintervall $C$ bleibt dabei unverändert.
Damit gilt $p\sim q$.

Ist $I=\varnothing$, so liegt einer der beiden Mediane an Position $1$. Für
den zugehörigen Wert gibt es keinen gleichartigen Eintrag links vom Median;
folglich gibt es auch keinen rechts davon. Dieser Wert tritt insgesamt genau
einmal auf. Alle Einträge in $J$ sind daher zwangsläufig gleich dem anderen
Wert, sodass aus $M(p)=M(q)$ bereits $p=q$ folgt. Der Fall
$J=\varnothing$ ist symmetrisch.
\end{proof}

\begin{cor}
Es sei $(n,p,q)$ wie in Definition
\ref{defn:problembeschreibung-median} gegeben. Genau dann ist $(n,p,q)$
lösbar, wenn die Anzahl der Nullen und die Anzahl der Einsen in $q$ ungerade
sind und
\[
M(p)=M(q)
\]
gilt.
\label{cor:algorithmus-kriterium-original}
\end{cor}

\begin{proof}
Nach Proposition~\ref{prop:paritaet-median} sind ungerade Anzahlen in $q$
notwendig. Sind sie erfüllt, folgt die Behauptung unmittelbar aus
Satz~\ref{thm:mediansignatur-charakterisierung}.
\end{proof}

\begin{cor}
Es lässt sich in $O(n)$ zeitlichem Aufwand bestimmen, ob das Problem lösbar
ist.
\end{cor}

\begin{proof}
Wir bestimmen in $p$ und $q$ jeweils die geordneten Positionen der Nullen und
Einsen. Das mittlere Element jeder Liste liefert $m_0$ beziehungsweise $m_1$.
Anschließend vergleichen wir $\ell$, $r$ und die beiden Teilwörter zwischen
diesen Grenzen. Jede Position wird nur konstant oft betrachtet. Daher beträgt
der Zeitaufwand $O(n)$. Werden die Positionen wie in der Python-Implementierung
gespeichert, beträgt der zusätzliche Speicherbedarf $O(n)$.
\end{proof}

\end{document}
